% --- обозначения ---
\newcommand{\ov}[1]{\overline{#1}}
\renewcommand{\emptyset}{\varnothing}
\newcommand{\R}{\mathbb R}
\newcommand{\N}{\mathbb N}
\newcommand{\Z}{\mathbb Z}
\newcommand{\Q}{\mathbb Q}
\newcommand{\M}{\mathfrak M}
\newcommand{\maj}{\operatorname{maj}}
\newcommand{\sd}{\mathbin{\triangle}}
\newcommand{\cupl}{\mathbin{\cup}}
% номер страницы оригинала на полях
\newcommand{\bookpage}[1]{\leavevmode\marginnote{\footnotesize\textcolor{gray}{[с.~#1]}}\ignorespaces}
% подпись к рисункам в стиле книги
\newcommand{\figcap}[1]{\par\smallskip\centering\small #1\par}
% --- добавленный материал ---
\definecolor{solcol}{RGB}{20,80,150}
\newtcolorbox{solution}[1][]{breakable,enhanced,colback=solcol!4,colframe=solcol!60,
  boxrule=0.5pt,arc=1pt,left=6pt,right=6pt,top=4pt,bottom=4pt,
  title={Решение упражнения #1},fonttitle=\bfseries\small,coltitle=white,colbacktitle=solcol!75}
\newtcolorbox{remarkbox}[1][]{breakable,enhanced,colback=orange!5,colframe=orange!60!black,
  boxrule=0.4pt,arc=1pt,title={#1},fonttitle=\bfseries\small}
% исправление опечатки: отметка в тексте
\newcommand{\fixed}[1]{\footnote{\textit{Исправлено:} #1}}
\theoremstyle{plain}
\newtheorem*{lemma*}{Лемма}
\newtheorem*{theorem*}{Теорема}
\newtheorem*{prop*}{Утверждение}
\newcommand{\exitem}[1]{\par\smallskip\noindent\textbf{#1}~}
